% problem-type: nhiet-cong-thuc-poisson
% order: 50
% Bo cuc: De vi du -> Cong thuc -> De + Bai giai

\section{Nhiệt toàn không gian --- công thức Poisson}

% ---------------------------------------------------------------
\subsection*{Đề bài ví dụ}
% ---------------------------------------------------------------
Cho $u$ là nghiệm bị chặn của
\[
  u_t=u_{xx}+f(x,t),\quad x\in\mathbb{R},\ t>0,\qquad u(x,0)=g(x).
\]
(i) Với $f=0$, $g=\chi_{[a,b]}$: chứng minh $u(x,t)\in(0,1)$. \quad
(ii) Với $f=e^{-t}\cos(cx)$, $g=\chi_{[a,b]}$: tính tường minh $u$.

\emph{\small(Nguồn: Đề thi MAT2306 --- PTĐHR 1, HK2 2024--2025, K67, Câu 4 (đề gốc $g=\chi_{[a-12,b]}$). Có đáp án chính thức.)}

% ---------------------------------------------------------------
\subsection*{Nhận ra dạng này khi nào?}
% ---------------------------------------------------------------
\begin{itemize}
  \item PDE \textbf{parabolic} $u_t=k\Delta u\,(+f)$ trên toàn không gian, cho dữ liệu ban đầu.
  \item Toàn $\mathbb{R}^n$, KHÔNG biên không gian. (Có biên $\Rightarrow$ xem dạng bài ``Nhiệt trên miền có biên''.)
\end{itemize}
\begin{meobox}
Dữ liệu là hàm đặc trưng $\Rightarrow$ đáp số ra $\operatorname{erf}$. Dữ liệu $\sin/\cos\Rightarrow$ dùng tích phân cos-Gauss (Phụ lục B).
\end{meobox}

% ---------------------------------------------------------------
\subsection*{Công thức \& phương pháp}
% ---------------------------------------------------------------
\subsubsection*{Công thức Poisson (nhân nhiệt)}
% method: cong-thuc-poisson-nhiet
\begin{dieukienbox}
Dùng khi: phương trình \textbf{nhiệt} $u_t=k\Delta u$ trên toàn $\mathbb{R}^n$, dữ liệu $u(\cdot,0)=g$.
\end{dieukienbox}

\begin{nhobox}
\textbf{Công thức Poisson} (tích chập với nhân nhiệt $\Phi$):
\[
  u(x,t)=\int_{\mathbb{R}^n}\Phi(x-y,t)\,g(y)\,dy,\qquad
  \Phi(x,t)=\frac{1}{(4\pi k t)^{n/2}}\,e^{-|x|^2/(4kt)}.
\]
1D: $\displaystyle u(x,t)=\frac{1}{\sqrt{4\pi k t}}\int_{-\infty}^{\infty} e^{-(x-y)^2/(4kt)}\,g(y)\,dy.$
\end{nhobox}

\begin{meobox}
$g$ là hàm đặc trưng $\chi_{[a,b]}$ $\Rightarrow$ đổi biến $w=\frac{y-x}{\sqrt{4kt}}$, tích phân ra
hàm $\operatorname{erf}$ (xem Phụ lục B). Nhiều chiều: nhân nhiệt tách tích theo từng biến.
\end{meobox}

\subsubsection*{Duhamel (khi có nguồn $f$)}
% method: duhamel-nhiet
\begin{dieukienbox}
Dùng khi: nhiệt \textbf{có nguồn} $u_t-k\Delta u=f$, dữ liệu ban đầu $=0$ (tách phần dữ liệu bằng chồng chất).
\end{dieukienbox}

\begin{nhobox}
\textbf{Duhamel cho nhiệt:} $\displaystyle u(x,t)=\int_0^t w(x,t;s)\,ds$, trong đó $w(\cdot,\cdot;s)$ giải
phương trình \emph{không nguồn} $w_t=k\Delta w$ ($t>s$) với \textbf{giá trị đầu} $w(\cdot,s;s)=f(\cdot,s)$.
\end{nhobox}

\begin{luuybox}
Khác sóng: với nhiệt (bậc 1 theo $t$), lát nguồn thành \emph{giá trị ban đầu} $w(\cdot,s;s)=f(\cdot,s)$;
với sóng (bậc 2) nó thành \emph{vận tốc đầu}. Tìm $w$ bằng công thức Poisson nhiệt với $g=f(\cdot,s)$, thay $t\to t-s$.
\end{luuybox}
\emph{(Khử điều kiện biên bằng thác triển: xem Phụ lục C.)}

% ---------------------------------------------------------------
\subsection*{Đề bài \& bài giải}
% ---------------------------------------------------------------
\paragraph{Nhắc lại đề.} $u_t=u_{xx}+f$ trên $\mathbb{R}$, $u(x,0)=g$. (i) $f=0,\ g=\chi_{[a,b]}$: chứng minh
$u\in(0,1)$. (ii) $f=e^{-t}\cos(cx),\ g=\chi_{[a,b]}$: tính tường minh.

\subsubsection*{(i) $f=0$, $g=\chi_{[a,b]}$}

\paragraph{Bước 1 --- ráp vào công thức Poisson.} Bài là nhiệt thuần nhất ($f=0$) trên $\mathbb{R}$ với dữ liệu đầu $g$. Áp công thức Poisson 1D (xem mục ``Công thức \& phương pháp'') với hệ số khuếch tán $k=1$ (vì PDE là $u_t=u_{xx}$):
\[
  u(x,t)=\frac{1}{\sqrt{4\pi t}}\int_{-\infty}^{\infty} e^{-(x-y)^2/(4t)}\,g(y)\,dy.
\]
Vì $g=\chi_{[a,b]}$ ($=1$ trên $[a,b]$, $=0$ ngoài) nên tích phân chỉ còn trên $[a,b]$:
\[
  u(x,t)=\frac{1}{\sqrt{4\pi t}}\int_a^b e^{-(x-y)^2/(4t)}\,dy.
\]

\paragraph{Bước 2 --- đổi biến để ra dạng Gauss chuẩn.} Mục tiêu: đưa mũ $-(x-y)^2/(4t)$ về $-z^2$ chuẩn. Đặt
\[
  z=\frac{y-x}{\sqrt{4t}}\ \Longrightarrow\ y=x+\sqrt{4t}\,z,\quad dy=\sqrt{4t}\,dz,\quad (x-y)^2/(4t)=z^2.
\]
Đổi cận theo $y$: $y=a\Rightarrow z=\tfrac{a-x}{\sqrt{4t}}$; $y=b\Rightarrow z=\tfrac{b-x}{\sqrt{4t}}$. Thay vào tích phân:
\[
  u=\frac{1}{\sqrt{4\pi t}}\cdot\sqrt{4t}\int_{(a-x)/\sqrt{4t}}^{(b-x)/\sqrt{4t}} e^{-z^2}\,dz
  =\frac{1}{\sqrt\pi}\int_{(a-x)/\sqrt{4t}}^{(b-x)/\sqrt{4t}} e^{-z^2}\,dz,
\]
trong đó $\dfrac{\sqrt{4t}}{\sqrt{4\pi t}}=\dfrac{1}{\sqrt\pi}$.

\paragraph{Bước 3 --- gọi tên bằng erf.} \textbf{Định nghĩa} hàm sai số:
\[
  \operatorname{erf}(w):=\frac{2}{\sqrt\pi}\int_0^w e^{-z^2}\,dz.
\]
Đây là một hàm lẻ, tăng nghiêm ngặt và $\operatorname{erf}(\pm\infty)=\pm1$, nên $\operatorname{erf}(w)\in(-1,1)$ với $w\in\mathbb{R}$. Từ định nghĩa, tách tích phân tại $0$:
\[
  \int_\alpha^\beta e^{-z^2}\,dz=\int_0^\beta e^{-z^2}dz-\int_0^\alpha e^{-z^2}dz=\frac{\sqrt\pi}{2}\big(\operatorname{erf}\beta-\operatorname{erf}\alpha\big).
\]
Áp vào kết quả Bước 2 với $\alpha=\tfrac{a-x}{\sqrt{4t}}$, $\beta=\tfrac{b-x}{\sqrt{4t}}$:
\[
  u(x,t)=\frac{1}{\sqrt\pi}\cdot\frac{\sqrt\pi}{2}\big(\operatorname{erf}\beta-\operatorname{erf}\alpha\big)
  =\tfrac12\Big(\operatorname{erf}\tfrac{b-x}{\sqrt{4t}}-\operatorname{erf}\tfrac{a-x}{\sqrt{4t}}\Big).
\]
\begin{nhobox}
Vì $b>a$ và $\operatorname{erf}$ tăng nên hiệu hai erf $>0$; lại $\operatorname{erf}\in(-1,1)$ nên hiệu $<2$.
Vậy $0<u(x,t)<1$.
\end{nhobox}

\subsubsection*{(ii) Thêm nguồn $f=e^{-t}\cos(cx)$}

\paragraph{Bước 1 --- chồng chất (tách thành 2 bài con).} Phương trình tuyến tính nên bài gốc
$u_t=u_{xx}+f,\ u(x,0)=g$ tách thành $u=u_g+u_f$, mỗi phần giải 1 bài đơn giản hơn:
\[
  (A)\ \begin{cases}(u_g)_t=(u_g)_{xx}\\ u_g(x,0)=g\end{cases}
  \qquad
  (B)\ \begin{cases}(u_f)_t=(u_f)_{xx}+f\\ u_f(x,0)=0\end{cases}
\]
Kiểm tra: cộng (A)$+$(B) cho $(u_g+u_f)_t=(u_g+u_f)_{xx}+f$ và dữ liệu đầu $g+0=g$ --- đúng bài gốc.
$u_g$ chính là phần (i) đã giải; còn $u_f$ (không dữ liệu đầu, có nguồn) giải bằng Duhamel dưới đây.

\paragraph{Bước 2 --- dựng sóng con $w$ (nguyên lý Duhamel).} \textbf{Nguyên lý Duhamel cho nhiệt} (xem mục ``Công thức \& phương pháp''): nghiệm của $u_{f,t}=u_{f,xx}+f$ với dữ liệu đầu $0$ là tổng đóng góp của các ``lát nguồn'' phát tại từng thời điểm $s\in(0,t)$:
\[
  u_f(x,t)=\int_0^t w(x,t;s)\,ds,
\]
trong đó với mỗi $s$ cố định, \emph{sóng con} $w(x,\tau;s)$ giải bài toán nhiệt \textbf{không nguồn} theo thời gian $\tau>s$:
\[
  w_\tau=w_{xx}\quad(\tau>s),\qquad w(x,s;s)=f(x,s)=e^{-s}\cos(cx).
\]
(Tức là nguồn ở thời điểm $s$ ``chích'' một giá trị đầu mới $f(\cdot,s)$ cho sóng con, rồi sóng con tự khuếch tán tự do từ $\tau=s$ tới $\tau=t$.) Mục tiêu Bước 3 dưới đây là tìm công thức tường minh cho $w$.

\paragraph{Bước 3 --- giải $w$ bằng hàm riêng.}\leavevmode

Sóng con $w(x,\tau;s)$ thoả bài toán nhiệt 1D \emph{không nguồn} theo biến thời gian $\tau$ (với $s$ là tham số cố định, $\tau$ chạy từ $s$ tới $t$):
\[
  w_\tau=w_{xx}\ (x\in\mathbb{R},\ \tau>s),\qquad w(x,s;s)=e^{-s}\cos(cx).
\]

\emph{Bước 3a --- $\cos(cx)$ là ``hàm riêng'' của toán tử $\partial_{xx}$.} Hàm $\varphi(x)$ gọi là \textbf{hàm riêng} của một toán tử tuyến tính $L$ nếu $L\varphi=\lambda\varphi$ với hằng số $\lambda$ (gọi là \emph{trị riêng}) --- nghĩa là tác động của $L$ chỉ là nhân với một con số, không trộn hàm. Ở đây, đạo hàm hai lần $\cos(cx)$:
\[
  \partial_{xx}\cos(cx)=-c^2\cos(cx),
\]
nên $\cos(cx)$ là hàm riêng của $\partial_{xx}$ với trị riêng $\lambda=-c^2$.

\emph{Bước 3b --- thử nghiệm dạng tách biến $w=A(\tau)\cos(cx)$.} Vì giá trị đầu của $w$ tại $\tau=s$ đã có dạng (hàm chỉ-theo-$\tau$) $\times\cos(cx)$, ta dự đoán $w$ giữ nguyên dạng đó với mọi $\tau$:
\[
  w(x,\tau;s)=A(\tau)\,\cos(cx),\qquad A\ \text{là hàm cần tìm}.
\]
Thay vào $w_\tau=w_{xx}$ và dùng $\partial_{xx}\cos(cx)=-c^2\cos(cx)$:
\[
  A'(\tau)\,\cos(cx)=A(\tau)\,\big(-c^2\cos(cx)\big)\ \Longrightarrow\ A'(\tau)=-c^2\,A(\tau).
\]
Phần $\cos(cx)$ \emph{triệt} hai bên (nó là hàm riêng), chỉ còn ODE thường cho $A$.

\emph{Bước 3c --- giải ODE cho $A$ với điều kiện đầu tại $\tau=s$.} Nghiệm tổng quát của $A'=-c^2 A$ là $A(\tau)=C\,e^{-c^2(\tau-s)}$ (chọn mốc tại $\tau=s$ cho gọn). Điều kiện đầu: tại $\tau=s$ ta cần $w(x,s;s)=e^{-s}\cos(cx)$, tức $A(s)=e^{-s}$. Vậy $C=e^{-s}$ và
\[
  A(\tau)=e^{-s}\,e^{-c^2(\tau-s)}.
\]

\emph{Bước 3d --- ráp lại tại $\tau=t$:}
\[
  w(x,t;s)=A(t)\cos(cx)=e^{-s}\,e^{-c^2(t-s)}\,\cos(cx).
\]

\paragraph{Bước 4 --- cộng theo $s$ (tính tích phân Duhamel).} Thay $w$ ở Bước 3d vào $u_f=\int_0^t w(x,t;s)\,ds$:
\[
  u_f(x,t)=\int_0^t e^{-s}\,e^{-c^2(t-s)}\,\cos(cx)\,ds.
\]
Giả sử $c^2\ne1$ (trường hợp $c=\pm1$ sẽ lấy giới hạn $L'$Hôpital ở cuối). Tách $\cos(cx)$ (không phụ thuộc $s$) ra ngoài và \emph{gom mũ} theo $s$:
\[
  -s-c^2(t-s)=-c^2 t+(c^2-1)s
  \ \Longrightarrow\ e^{-s}\,e^{-c^2(t-s)}=e^{-c^2 t}\,e^{(c^2-1)s}.
\]
Vậy
\[
  u_f(x,t)=\cos(cx)\,e^{-c^2 t}\int_0^t e^{(c^2-1)s}\,ds.
\]
Tích phân mũ cơ bản (nguyên hàm $\frac{1}{c^2-1}e^{(c^2-1)s}$):
\[
  \int_0^t e^{(c^2-1)s}\,ds=\frac{e^{(c^2-1)t}-1}{c^2-1}.
\]
Thay lại và phân phối $e^{-c^2 t}$ vào tử:
\[
  u_f(x,t)=\cos(cx)\,e^{-c^2 t}\cdot\frac{e^{(c^2-1)t}-1}{c^2-1}
  =\cos(cx)\cdot\frac{e^{-c^2 t}\,e^{(c^2-1)t}-e^{-c^2 t}}{c^2-1}.
\]
Mũ đầu rút gọn: $-c^2 t+(c^2-1)t=-t$, nên
\[
  u_f(x,t)=\frac{e^{-t}-e^{-c^2 t}}{c^2-1}\,\cos(cx).
\]
(Trường hợp $c^2=1$: lấy giới hạn cho $u_f=t\,e^{-t}\cos(cx)$.) Cuối cùng $u=u_g+u_f$.
\begin{nhobox}
\[
  u(x,t)=\tfrac12\Big(\operatorname{erf}\tfrac{b-x}{\sqrt{4t}}-\operatorname{erf}\tfrac{a-x}{\sqrt{4t}}\Big)
  +\frac{e^{-t}-e^{-c^2 t}}{c^2-1}\cos(cx).
\]
\end{nhobox}